\documentclass[a4paper,12pt]{article}
\usepackage[T2A]{fontenc}
\usepackage[utf8]{inputenc}
\usepackage[russian]{babel}
\usepackage{amsmath,amssymb,icomma}
\usepackage[a4paper,margin=23mm,headheight=15pt,footskip=15mm]{geometry}
\usepackage{microtype,booktabs,tabularx,array,xcolor,tikz,fancyhdr}
\usepackage[colorlinks=true,linkcolor=teal!70!black]{hyperref}
\definecolor{accent}{HTML}{008080}
\definecolor{darkaccent}{HTML}{004D4D}
\definecolor{error}{HTML}{993B37}
\setlength{\parindent}{0pt}
\setlength{\parskip}{5pt}
\renewcommand{\arraystretch}{1.5}
\pagestyle{fancy}
\fancyhf{}
\fancyhead[L]{\small Теория вероятностей}
\fancyhead[R]{\small Марк Гукин}
\fancyfoot[L]{\footnotesize Лёвин Артём Александрович эксклюзивно для Марка Гукина}
\fancyfoot[R]{\small\thepage}
\renewcommand{\headrulewidth}{0.3pt}
\newcommand{\topic}[1]{\par\medskip\textcolor{darkaccent}{\textbf{#1}}\par\nopagebreak}
\newcommand{\err}[1]{\textcolor{error}{\textbf{Ошибка:}} #1}
\newcommand{\task}[2]{\phantomsection\label{task:#1}{\Large\color{accent} #1. #2}\par\medskip}
\begin{document}
\begin{titlepage}
\thispagestyle{empty}
\centering
\vspace*{35mm}
{\large\color{darkaccent} МАТЕМАТИКА\par}
\vspace{18mm}
{\LARGE Ревью домашней работы\par}
\vspace{12mm}
{\Large\color{accent} Теория вероятностей\par}
\vfill
{\large Марк Гукин\par}
\vspace{10mm}
25 ноября 2026 г.\par
\vspace{16mm}
{\small Лёвин Артём Александрович\par эксклюзивно для Марка Гукина\par}
\vspace{18mm}
\begin{tikzpicture}
\draw[accent,line width=0.6pt] (0,0) .. controls (2,0.7) and (4,-0.7) .. (6,0) .. controls (8,0.7) and (10,-0.7) .. (12,0);
\end{tikzpicture}
\vspace{5mm}
\end{titlepage}
\section*{Сводная таблица}
В таблице приведены задания, требующие коррекции. Номер задания ведёт к подробному разбору.
\par\medskip
\begin{tabularx}{\linewidth}{@{}c >{\raggedright\arraybackslash}p{38mm} >{\raggedright\arraybackslash}X p{22mm} p{23mm}@{}}
\toprule
№ & Тема & Суть ошибки & \shortstack[l]{Ваш\\ответ} & \shortstack[l]{Верный\\ответ}\\
\midrule
\hyperref[task:4]{4} & Противоположное событие & Ошибка при вычитании десятичной дроби из единицы. & $0,9735$ & $0,9375$\\
\hyperref[task:7]{7} & Полная вероятность & Во второй ветви выбрана вероятность брака вместо вероятности исправности. & $0,5615$* & $0,9665$\\
\bottomrule
\end{tabularx}\par
\smallskip
{\small * Значение восстановлено однозначно по записанной сумме произведений; в рукописи в конце читаются цифры 5615.}
\clearpage
\thispagestyle{empty}
\vspace*{\fill}
\begin{center}
{\Huge\color{accent} Разбор ошибок}
\end{center}
\vspace*{\fill}
\clearpage
\task{4}{Четыре последние цифры номера}
\topic{Условие}
Каждая из четырёх позиций номера независимо заполняется равновероятной цифрой от 0 до 9. Нуль на первой позиции допустим. Найдите вероятность того, что произведение четырёх цифр будет чётным.
\topic{Ваше решение и ответ}
Вы правильно перешли к противоположному событию и получили
\[
\left(\frac{5}{10}\right)^4=0,0625;\qquad 1-0,0625=0,9735.
\]
\err{Арифметическая неточность в последнем равенстве.} Последний верный результат --- $0,0625$. При вычитании из единицы получено неверное число. Метод решения сохраняем полностью.
\topic{Ключевая идея}
Произведение целых чисел чётно, если хотя бы один множитель чётный. Нуль тоже является чётным числом. Противоположное событие --- все четыре цифры нечётные. Эти два события несовместны и вместе охватывают все исходы, поэтому сумма их вероятностей равна единице.
\topic{Полное правильное решение}
Нечётных цифр пять: $1,3,5,7,9$. Вероятность нечётной цифры на любой позиции равна $5/10=0,5$. Позиции независимы, поэтому вероятность четырёх нечётных цифр равна
\[
P(\overline A)=0,5^4=0,25\cdot0,25=0,0625.
\]
Здесь $A$ --- событие «произведение чётное». По формуле противоположного события
\[
P(A)=1-P(\overline A)=1,0000-0,0625=0,9375.
\]
При вычитании удобно дописать нули после запятой и вычитать по разрядам. Все случаи с одной, двумя, тремя и четырьмя чётными цифрами уже учтены.
\topic{Проверка и итог}
$0,9375+0,0625=1$. Дополнительно: из $10^4=10000$ последовательностей ровно $5^4=625$ состоят из нечётных цифр. Остаётся $9375$ подходящих последовательностей. Ответ: \textbf{$0,9375$}.
\topic{Задача на закрепление}
Код из пяти символов независимо заполняется цифрами от 0 до 9 с одинаковой вероятностью. Найдите вероятность того, что произведение всех пяти цифр будет чётным. Выполните вычисления точно.
\clearpage
\task{7}{Поставка электрических ламп}
\topic{Условие}
Первый завод поставляет 55\% ламп, второй --- 45\%. Доля бракованных ламп на первом заводе составляет 2\%, на втором --- 5\%. Найдите вероятность того, что случайно выбранная из общей партии лампа исправна. Отделите вероятность выбора завода от вероятности качества лампы.
\topic{Ваше решение и ответ}
Вы верно построили дерево и нашли вероятности исправности на обоих заводах: $0,98$ и $0,95$. В итоговой сумме записано
\[
0,55\cdot0,98+0,45\cdot0,05=0,5615.
\]
\err{Во втором произведении выбрана ветвь брака.} Первое произведение $0,539$ верно. Второе произведение $0,0225$ соответствует неисправной лампе второго завода. В итоге сложены вероятности двух исходов с разным качеством.
\topic{Ключевая идея}
Обозначим через $A$ событие «лампа исправна», через $H_1,H_2$ --- выбор лампы первого и второго заводов. Вдоль каждой ветви перемножаем вероятности, а затем складываем только пути, ведущие к одному и тому же искомому событию.
\begin{center}
\begin{tikzpicture}[x=1cm,y=1cm,every node/.style={font=\small},edge/.style={draw=accent,line width=0.7pt}]
\node (s) at (0,0) {Выбор лампы};
\node (h1) at (-3,-1.1) {I завод}; \node (h2) at (3,-1.1) {II завод};
\node (a1) at (-4.4,-2.5) {Исправна}; \node (b1) at (-1.6,-2.5) {Брак};
\node (a2) at (1.6,-2.5) {Исправна}; \node (b2) at (4.4,-2.5) {Брак};
\draw[edge] (s)--node[above left] {$0,55$}(h1);
\draw[edge] (s)--node[above right] {$0,45$}(h2);
\draw[edge] (h1)--node[left] {$0,98$}(a1);
\draw[gray,dashed] (h1)--node[right] {$0,02$}(b1);
\draw[edge] (h2)--node[left] {$0,95$}(a2);
\draw[gray,dashed] (h2)--node[right] {$0,05$}(b2);
\node[accent] at (-4.4,-3.05) {$0,539$}; \node[accent] at (1.6,-3.05) {$0,4275$};
\end{tikzpicture}
\end{center}
\topic{Исправление и полное решение}
Вероятности выбора заводов: $P(H_1)=0,55$, $P(H_2)=0,45$. Условные вероятности исправности: $P(A\mid H_1)=1-0,02=0,98$, $P(A\mid H_2)=1-0,05=0,95$. По формуле полной вероятности
\begin{align*}
P(A)&=P(H_1)P(A\mid H_1)+P(H_2)P(A\mid H_2)\\
&=0,55\cdot0,98+0,45\cdot0,95\\
&=0,539+0,4275=0,9665.
\end{align*}
\topic{Проверка и итог}
Через брак: $0,55\cdot0,02+0,45\cdot0,05=0,0335$. Значит, $1-0,0335=0,9665$. Кроме того, общий показатель исправности должен лежать между $0,95$ и $0,98$. Ответ: \textbf{$0,9665$}.
\topic{Задача на закрепление}
Две линии выпускают 60\% и 40\% датчиков соответственно. Доля дефектных датчиков на первой линии равна 3\%, на второй --- 8\%. Найдите вероятность исправности случайно выбранного датчика из общей партии. Изобразите дерево вероятностей.
\end{document}
